% 来源：docs/03-autoregressive-models.md
% 仅建立章节框架；正文、公式、图表与习题内容待后续迁移。
\chapter{自回归模型}
\label{chap:ar}

\section*{本章要点}
% TODO：迁移原章“先抓住这几件事”的要点。
\contentplaceholder

\section{先确定顺序，再建模条件分布}
\label{sec:ar-ordering}
% TODO：补充本节内容。
\contentplaceholder

\subsection{联合概率分解与归一化}
\label{subsec:ar-chain-rule}
% TODO：补充本节内容。
\contentplaceholder

\subsection{变量顺序与归纳偏置}
\label{subsec:ar-inductive-bias}
% TODO：补充本节内容。
\contentplaceholder

\section{网络输出的是分布参数}
\label{sec:ar-outputs}
% TODO：补充本节内容。
\contentplaceholder

\subsection{二元变量与伯努利输出}
\label{subsec:ar-binary}
% TODO：补充本节内容。
\contentplaceholder

\subsection{类别变量与连续变量的输出}
\label{subsec:ar-categorical-continuous}
% TODO：补充本节内容。
\contentplaceholder

\subsection{概率参数、采样与确定性决策}
\label{subsec:ar-sampling-decisions}
% TODO：补充本节内容。
\contentplaceholder

\section{从简单模型到共享网络}
\label{sec:ar-shared-networks}
% TODO：补充本节内容。
\contentplaceholder

\subsection{FVSBN：用逻辑回归表示条件概率}
\label{subsec:ar-fvsbn}
% TODO：补充本节内容。
\contentplaceholder

\subsection{NADE：非线性条件分布与参数共享}
\label{subsec:ar-nade}
% TODO：补充本节内容。
\contentplaceholder

\subsection{MADE：在一个网络内施加掩码}
\label{subsec:ar-made}
% TODO：补充本节内容。
\contentplaceholder

\section{架构不同，因果约束相同}
\label{sec:ar-architectures}
% TODO：补充本节内容。
\contentplaceholder

\subsection{RNN 与 PixelRNN}
\label{subsec:ar-rnn}
% TODO：补充本节内容。
\contentplaceholder

\subsection{因果 CNN 与 WaveNet}
\label{subsec:ar-causal-cnn}
% TODO：补充本节内容。
\contentplaceholder

\subsection{PixelCNN 的空间与通道掩码}
\label{subsec:ar-pixelcnn}
% TODO：补充本节内容。
\contentplaceholder

\subsection{Transformer 的输入移位与因果注意力}
\label{subsec:ar-transformer}
% TODO：补充本节内容。
\contentplaceholder

\section{生成、似然计算与训练要分开看}
\label{sec:ar-computation}
% TODO：补充本节内容。
\contentplaceholder

\subsection{生成：祖先采样}
\label{subsec:ar-generation}
% TODO：补充本节内容。
\contentplaceholder

\subsection{似然计算：观测已经完整}
\label{subsec:ar-likelihood}
% TODO：补充本节内容。
\contentplaceholder

\subsection{训练：Teacher Forcing}
\label{subsec:ar-teacher-forcing}
% TODO：补充本节内容。
\contentplaceholder

\subsection{并行计算、曝光偏差与误差积累}
\label{subsec:ar-parallelism}
% TODO：补充本节内容。
\contentplaceholder

\section{这类模型适合什么？}
\label{sec:ar-strengths-limits}
% TODO：补充本节内容。
\contentplaceholder

\subsection{似然建模的优势与顺序生成的代价}
\label{subsec:ar-tradeoffs}
% TODO：补充本节内容。
\contentplaceholder

\subsection{连续数据与潜在 token 表示}
\label{subsec:ar-representations}
% TODO：补充本节内容。
\contentplaceholder

\section{自检}
\label{sec:ar-review}
% TODO：补充本节内容。
\contentplaceholder

\section{来源与延伸阅读}
\label{sec:ar-reading}
% TODO：补充本节内容。
\contentplaceholder
